\documentclass[11pt,a4paper]{article}

% ---------- Packages ----------
\usepackage[utf8]{inputenc}
\usepackage[T1]{fontenc}
\usepackage{lmodern}
\usepackage{graphicx}
\usepackage{amsmath, amssymb}
\usepackage{booktabs}
\usepackage{hyperref}
\usepackage[margin=1in]{geometry}
\usepackage{caption}
\usepackage{subcaption}
\usepackage{enumitem}
\usepackage{float}     % for [H] exact figure placement
\usepackage{placeins}  % for \FloatBarrier

\hypersetup{
    colorlinks=true,
    linkcolor=blue,
    citecolor=blue,
    urlcolor=blue
}

% Title 
\title{\textbf{Emotion Recognition from Text using BERT and Classical Machine Learning Models}}
\author{Md Naim Hassan Saykat}
\date{}

\begin{document}
\maketitle

\begin{abstract}
This project addresses the task of multi-class emotion recognition from text. We compare a transformer-based BERT model with classical machine learning approaches including Support Vector Machines (SVM), Random Forest (RF), and Logistic Regression (LR). We also evaluate a simple majority-vote ensemble over all models. Results on the \texttt{dair-ai/emotion} dataset are reported using test accuracy and macro-averaged precision, recall and F1. BERT achieves the strongest overall performance, while the ensemble provides a balanced alternative across classes.
\end{abstract}

\noindent \textbf{Keywords:} Emotion recognition, BERT, SVM, Random Forest, Logistic Regression, Ensemble, NLP.

% Introduction 
\section{Introduction}
Automatic emotion recognition from text is a key problem in natural language processing (NLP) with applications in social media analytics, mental-health monitoring and conversational AI. Unlike sentiment analysis, which is typically binary (positive/negative), emotion recognition is a multi-class problem that must distinguish among categories such as \textit{joy}, \textit{sadness}, \textit{anger}, \textit{fear}, \textit{love} and \textit{surprise}.

Recent advances in deep learning, particularly transformers such as BERT \cite{devlin2019bert}, have enabled substantial improvements in NLP classification tasks. In this work we use the \texttt{dair-ai/emotion} dataset \cite{saravia2018carer}. Classical baselines (SVM, Random Forest, Logistic Regression) are implemented with scikit-learn \cite{pedregosa2011scikit}. We compare these approaches and a simple ensemble to understand quality/compute trade-offs for the task.

% Methodology 
\section{Methodology}
\subsection{Dataset}
We use the Hugging Face \texttt{dair-ai/emotion} dataset \cite{saravia2018carer}, which contains English texts annotated with six emotion classes: \textit{sadness}, \textit{joy}, \textit{love}, \textit{anger}, \textit{fear} and \textit{surprise}. We follow the official train/validation/test split.

\subsection{Models}
\begin{itemize}[leftmargin=1.2em]
    \item \textbf{BERT:} \texttt{bert-base-uncased} fine-tuned for sequence classification with six output labels. Tokenization uses \texttt{AutoTokenizer}.
    \item \textbf{SVM:} Linear-kernel SVM trained on TF--IDF features.
    \item \textbf{Random Forest:} 100 trees trained on TF--IDF features.
    \item \textbf{Logistic Regression:} Linear baseline with TF--IDF features.
    \item \textbf{Ensemble:} Simple majority vote over \{BERT, SVM, RF, LR\}.
\end{itemize}

\subsection{Evaluation Metrics}
We report test accuracy and macro-averaged precision, recall and F1-score computed with scikit-learn.

% Experiments 
\section{Experiments}
\subsection{Training Setup}
BERT is fine-tuned for 3 epochs with AdamW and learning rate $2\times 10^{-5}$. Classical models use scikit-learn defaults (SVM with a linear kernel, RF with 100 estimators, LR with regularization as in \texttt{liblinear}). All results are measured on the held-out test split.

\subsection{Results}
Table~\ref{tab:results} summarizes the results.\footnote{Macro metrics are the macro-averaged values from the printed classification reports in the notebook. For Random Forest, the run that was logged only printed accuracy; macro metrics were not printed in that run and are thus left blank.}

\begin{table}[H]
\centering
\begin{tabular}{lcccc}
\toprule
\textbf{Model} & \textbf{Test Accuracy} & \textbf{Macro Precision} & \textbf{Macro Recall} & \textbf{Macro F1} \\
\midrule
BERT & 0.92 & 0.87 & 0.88 & 0.87 \\
SVM & 0.89 & 0.85 & 0.81 & 0.83 \\
Random Forest & 0.88 & N/A & N/A & N/A \\
Logistic Regression & 0.87 & 0.86 & 0.78 & 0.81 \\
Ensemble (BERT+SVM+RF+LR) & 0.90 & 0.90 & 0.80 & 0.83 \\
\bottomrule
\end{tabular}
\caption{Comparison of models on the \texttt{dair-ai/emotion} dataset.}
\label{tab:results}
\end{table}

\subsection{Qualitative observations}
Per-class reports show that \textit{surprise} is the hardest category across models (lower recall), while \textit{sadness} and \textit{joy} are recognized most reliably. The ensemble improves robustness for several minority classes by aggregating complementary errors.

% Discussion 
\section{Discussion}
BERT delivers the highest test accuracy and macro F1, as expected for a contextual transformer. However, SVM and LR remain strong baselines given their much lower computational cost. The majority-vote ensemble offers a simple way to balance precision and recall across classes without additional tuning.

% Conclusion 
\section{Conclusion}
We presented a comparative study of deep and classical approaches for multi-class emotion recognition. BERT provides the best overall performance, classical baselines offer efficient alternatives, and a simple ensemble yields competitive, balanced results. Future work includes computing and logging the full Random Forest classification report, exploring weighted voting or stacking, and evaluating larger transformer backbones (e.g., RoBERTa, DeBERTa).

% References 
\bibliographystyle{plain}
\bibliography{references}

\end{document}